\scanpage{227}
Quelques calculs d'homotopie

Soit $k$ un anneau de base comm.,
$\Phi_*=\Phi^k_{*\unsure{k}}$ le $k$-module simplicial $I\mathbin{-}k^I$,
$k_*$ le $k$-module simplicial cste défini par $k$, $I\mathbin{-}k$.
On a une application diagonale
\[k_*\mathbin{-}\Phi_*\]
d'où un quotient
\[\Psi_*=\Phi_*/k_*:I\longrightarrow k^I/k.\]

\uline{Proposition 1} $\Phi_*$ est homotopique à $0$ \deleted{$\unsure{\Phi_k}$}
(\deleted{module}\inserted{i.e. au $k$-}simplicial $0$).

Par la théorie de Puppe-Dold, il suffit de voir que le complexe de chaines
\inserted{$\Phi_*^!$} associé \deleted{$\unsure{\Phi}$} à $\Phi$ est hom. à $0$.
Pour ceci, on se rappelle que
\[
\Phi_n^!=\{x\in\Phi_n\mid\partial_1^*(x)=- -=\partial_n^*(x)=0\}
\]
d'où aussitôt
\[
\Phi_n^!=\begin{cases}0&\text{si }n\geq2\\k&\text{si }n=0,1\end{cases}
\]
l'op. \deleted{\unsure{cob}} bord
\diagram{\begin{tikzcd}[ampersand replacement=\&]
\partial=\partial_0:\&\Phi_1^!\arrow[r]\arrow[d,"\sim"{sloped},phantom]\&\Phi_0^!\arrow[d,"\sim"{sloped},phantom]\\
\&k\&k
\end{tikzcd}}
étant l'identité, d'où la conclusion.

\uline{Corollaire 1.1} Pour tt. foncteur $F:\operatorname{Mod}(k)\longrightarrow C$
($C$ une \unsure{catégorie}) $F(\Phi_*)$ est homotopique à $F(0)$.
En particulier
\[
\left\{\begin{aligned}
\Lambda_k^i\Phi_*\ \inserted{(\Gamma_k^i\Phi_*,\operatorname{Sym}_k^i\Phi_*)}
&\text{ est homotopique à }\Lambda_k^i(0)\ \inserted{(\Gamma_k^i(0),\operatorname{Sym}_k^i(0))}
=\begin{cases}0&\text{si }i\ne0\\(k)_*&\text{si }i=0\end{cases}\\
\deleted{\Gamma_k^i\Phi_*}&\deleted{\mathbin{-}\Gamma_k^i(0)=\begin{cases}0&\text{si }i\ne0\\k_*&\text{si }i=0\end{cases}}
\end{aligned}\right.
\]
\margin{\figure{\begin{tikzpicture}\draw (0,2.4)--(0,0)--(.12,0);\end{tikzpicture}}}
\uline{Corollaire 1.2.} Les groupes d'homotopie de $\Psi_*=\Phi_*/k_*$ sont
\[
\pi_i(\deleted{\Phi_*}\inserted{\Psi_*}\deleted{/k_*})
\simeq\pi_{i-1}(k_*)\simeq\begin{cases}0&\text{si }i\ne1\\k&\text{si }i=1\end{cases}.
\]
En d'autres termes $\deleted{\Phi_*}\inserted{\Psi_*}\deleted{/k_*}$ est un $K_*(k,1)$.
[En fait, on a $(\unsure{\Psi|\Phi}_*^!\deleted{/k_*^!})=$ complexe réduit à $k$ \uline{placé en dimension $1$}.

On conclut de ceci \deleted{\unsure{et des props. de Dold-Puppe}}

\uline{Corollaire 1.3.}
\[
\pi_j(\Lambda^i(\Phi_*/k_*))\simeq
\begin{cases}0&\text{si }j\ne i\\k&\text{si }j=i\end{cases}
\]
(i.e. $\Lambda^i\Phi_*/k_*$ est un $K_*(k,i)$).
\margin{226}
